\section{关键词演化与热点趋势}

\begin{reportbox}{分析目标}
关键词和主题演化用于识别近五年科技文献中的热点变化。本章不只展示高频词, 还关注年度归一化占比、近期动量和主题模型是否给出一致证据, 以区分真实升温和样本规模变化造成的表观增长。
\end{reportbox}

\subsection{处理口径}

本章先对大小写、连字符、复数和常见同义表达进行归并, 例如 retrieval-augmented generation、RAG 与 retrieval augmented generation 视为同一方向; 再融合显式关键词和题名摘要中的受控术语信号, 过滤 study、method、model 等低信息量泛词。每个主题词同时计算融合文档数、显式关键词支持数、文本术语支持数、年度占比、近期动量、共现网络位置、burst window、生命周期阶段和语义漂移代理指标。

从整体结果看, retrieval augmented generation、large language model、prompt engineering、vision language model、knowledge graph 等方向同时具有较高频次和较强近期动量。它们不仅出现在词频统计中, 也能在年度归一化趋势和主题模型中得到支持, 因而比单纯高频词更适合作为热点候选。作者边和关键词边均使用增量缓存保存单篇论文的边贡献; 数据继续扩充时, 未变化论文可复用已有边贡献, 只重算新增或内容变化的论文。

\begin{methodnote}
2026 年需要单独解释: 截至 2026 年 6 月 21 日, 该年度尚未结束, 且新近论文的关键词和索引记录存在更新滞后。报告主图已使用显式关键词与题名摘要术语的融合信号进行修正, 2026 年融合信号比单看显式关键词增加约 22.7\%; 因此 2026 年只作为未完整年度的早期信号, 不与完整年度作机械对比。
\end{methodnote}

\subsection{高频与升温主题}
本节按"频次 -> 年度热度 -> 近期动量 -> 共现结构 -> 生命周期/主题模型"的顺序阅读。图~\ref{fig:top-keywords} 到图~\ref{fig:keyword-trends} 给出主趋势,图~\ref{fig:keyword-network} 以后用于解释结构与边界。

\begin{figure}[H]
  \centering
  \figasset[0.92\textwidth]{top_keywords.png}
  \caption{2022--2026 年高频融合主题词。数据来源:\texttt{data/analysis/keyword\_trends.csv};统计口径为显式关键词与题名摘要术语的融合信号。}
  \label{fig:top-keywords}
\end{figure}

\begin{figure}[!htbp]
  \centering
  \figasset[0.94\textwidth]{keyword_evolution_heatmap.png}
  \caption{融合主题词年度热度矩阵。数据来源:\texttt{data/analysis/keyword\_trends.csv};每行按自身最大值归一化,2026 年仅作为早期信号。}
  \label{fig:keyword-heatmap}
\end{figure}

\begin{figure}[!htbp]
  \centering
  \figasset[0.9\textwidth]{keyword_momentum.png}
  \caption{融合主题词近期动量。数据来源:\texttt{data/analysis/keyword\_trends.csv};该指标比较 2024--2025 与 2022--2023 的相对变化,不把 2026 年未完整数据纳入完整年度动量。}
  \label{fig:keyword-momentum}
\end{figure}

\begin{figure}[!htbp]
  \centering
  \figasset[0.92\textwidth]{keyword_normalized_trends.png}
  \caption{代表性融合主题词年度归一化文档占比。数据来源:\texttt{data/analysis/keyword\_trends.csv};归一化后仍持续上升的主题更可能反映真实研究关注变化。}
  \label{fig:keyword-trends}
\end{figure}

\begin{figure}[!htbp]
  \centering
  \figasset[\textwidth]{keyword_cooccurrence_network.png}
  \caption{关键词共现网络。数据来源:\texttt{data/analysis/keyword\_network.*};节点大小表示共现加权度,节点颜色表示关键词社区,边宽表示同篇论文中的分数化共现强度。该图用于识别 RAG、knowledge graph、question answering、clinical search 等主题之间的结构连接。}
  \label{fig:keyword-network}
\end{figure}

\begin{figure}[!htbp]
  \centering
  \figasset[0.92\textwidth]{keyword_lifecycle.png}
  \caption{关键词生命周期阶段。横线表示首次出现、峰值年份和最近活跃年份的区间, 圆点大小表示融合文档支持度, 阶段分为 emergence、growth、maturity、decline。}
\end{figure}

\begin{figure}[!htbp]
  \centering
  \figasset[0.92\textwidth]{keyword_burst_windows.png}
  \caption{关键词年度突发窗口。该图标出年度增长、萌发或回落较强的主题词, 用于把“长期高频”与“近期突然变化”区分开。}
\end{figure}

\begin{figure}[!htbp]
  \centering
  \figasset[0.92\textwidth]{topic_year_share.png}
  \caption{主题模型年度占比。主题模型用于检查热点词是否对应稳定的语义结构, 而不是孤立词频波动。}
\end{figure}

\FloatBarrier

\subsection{热点分层}

\begin{tableblock}
\begin{tabularx}{\textwidth}{>{\bfseries}p{31mm}X X}
\toprule
候选方向 & 数据信号 & 解释边界 \\
\midrule
RAG & 融合信号覆盖 5,579 篇候选文档, 其中 628 篇具有显式关键词支持, 题名摘要信号在近年明显增强 & 强上升信号, 但需要注意计算机和医学方向样本共同推动。 \\
LLM & 融合信号覆盖 17,600 篇候选文档, 其中 2,081 篇具有显式关键词支持, 与 prompt engineering、multimodal learning 共同升温 & 稳定增长信号, 需区分通用模型热度和具体科研应用。 \\
Knowledge graph & 融合信号覆盖 8,288 篇候选文档, 显式关键词和题名摘要均提供支持, 作者合作社区中也能观察到相关团簇 & 稳定主题结构, 更适合解释知识组织和跨学科连接。 \\
Materials informatics & materials informatics、battery、catalyst discovery 等术语在材料方向具有持续信号, 但更多依赖题名摘要识别 & 样本规模低于计算机和生物医学, 应优先看归一化趋势。 \\
\bottomrule
\end{tabularx}
\end{tableblock}

\subsection{趋势判定口径}

本报告不把单一高频词直接写成热点, 而是先消除年度样本规模差异, 再观察近期动量与多证据一致性。设关键词 $k$ 在年份 $y$ 的融合命中文档数为 $n_{k,y}$, 当年分析文献总数为 $N_y$, 年度归一化热度定义为:
\formulaline{
h_{k,y}=\frac{n_{k,y}}{N_y}
}{其中 $h_{k,y}$ 用于比较不同年份的相对关注度, 避免文献总量增长造成表观升温。}

近期动量用于区分稳定高频与真正升温。报告主口径比较 2024--2025 与 2022--2023 的平均归一化热度:
\formulaline{
m_k=\frac{h_{k,2024}+h_{k,2025}}{2}-\frac{h_{k,2022}+h_{k,2023}}{2}
}{2026 年为未完整年度, 只作为早期信号进入解释, 不纳入完整年度动量。}

最终写入结论时, 本报告将主题分为“上升信号”“稳定增长”“早期萌芽”或“可能噪声”, 不使用确定性预测。只有当频次、年度占比、近期动量、共现结构、生命周期/突发窗口、主题模型和代表论文多个证据同时支持时, 一个方向才进入强上升候选; 若某个词只在共现网络中靠近核心主题, 但年度占比和突发窗口不支持, 则只作为关联主题而非热点结论。
